\section{Tres Product-Critical Systems: Indexing, Inference y Apply}

El curso primero divide el sistema en indexing, model inference y la product/apply layer. Los tres sirven al mismo editor, pero tienen curvas de recursos completamente distintas: indexing tiende a ser asíncrono y data-intensive; inference está limitado por la GPU capacity, el batching y la latency; la apply layer debe trasladar los resultados generados de forma rápida y segura al archivo que el usuario está editando.

\subsection{Cada workload requiere su propio SLO}

Esta sección convierte las diferencias de workload del párrafo anterior en un contrato observable. Al leer la figura, primero se define por separado qué significa «éxito, degradación y falla» para indexing, inference, apply y el data plane, y después se eligen métricas como freshness, recall, queue time, edit success o job correctness. Si todas las rutas comparten una sola cifra de availability, durante un incidente no es posible determinar qué tipo de trabajo conviene sacrificar para proteger la ruta principal de interacción.

\lecturefigure{02-three-systems.png}{La escala de Cursor proviene de cuatro tipos de carga: indexing, inference, product/apply y data plane.}{Subtítulos locales 02:20--05:20; redibujo conceptual.}

\begin{knowledgebox}{Lectura de la figura: primero, definir la falla de cada sistema}
En indexing, la falla puede ser un índice stale, missing o con recall bajo; en inference, puede ser un timeout, un rate limit o baja calidad; en apply, puede ser un conflicto de patch, una posición errónea o un editor freeze; en el data plane, puede ser un job perdido, una ejecución duplicada o un billing que no se puede explicar. Solo si primero se clasifican los tipos, las alertas y las degradaciones no se contaminan entre sí.
\end{knowledgebox}

\begin{importantbox}{Un SLO no puede limitarse a «99.9\% available»}
\begin{itemize}
  \item Index freshness: el retraso entre un workspace change y el momento en que es recuperable;
  \item Retrieval quality: recall/utility bajo una latency y un filter determinados;
  \item Inference: time-to-first-token, tokens/s, error y queue time;
  \item Apply: edit success, conflict, validation latency y undo safety.
\end{itemize}
\end{importantbox}

\subsection{Apply model: separar «qué se quiere cambiar» de «cómo se escribe en disco»}

El modelo de planificación es bueno para entender la intención, elegir archivos y generar modificaciones; el specialized apply model, en cambio, se encarga de la coincidencia de posiciones, la reescritura local y la token application de alta velocidad. Se separan porque el modelo más potente no necesariamente es el mejor en la patch application determinista y de baja latencia. El sistema necesita diseñar interfaces distintas para «pensar» y para «escribir en el editor state».

\teachervoice{En clase se insiste repetidamente en que la capa de producto no es una UI delgada colocada después de la salida del modelo. Lo que realmente hace que una herramienta «se sienta cómoda de usar» suele ser el trabajo de sistemas que difícilmente aparece en un benchmark: apply, streaming, context preparation y editor integration.}

\subsection{Resumen de la sección}

Una plataforma de AI coding no es un único servicio de backend, sino una combinación de varios sistemas con SLO y formas de datos distintos. El valor del producto surge de la coordinación de estos sistemas dentro de una misma interacción, no de la métrica aislada de un modelo.
